\subsection{Geoparsing and Named Entity Recognition}

The extraction of geographic information from textual sources commonly relies on geoparsing pipelines that combine named entity recognition (NER) with geocoding techniques \cite{Gritta2018, Karimzadeh2019Geoparsing}. Within environmental and disaster-related applications, these approaches are frequently used to identify locations associated with reported impacts in news articles, reports, or social media content.

Most geoparsing systems first detect location mentions using NER models and subsequently map these entities to geographic coordinates or administrative regions through gazetteers such as GeoNames \cite{Gritta2018}. Traditional approaches often employ rule-based systems or statistical NLP models, while more recent methods increasingly rely on transformer-based architectures for contextual entity recognition and disambiguation.

In drought impact research, Sodoge et al.~\cite{sodoge2023} combine spaCy-based NER with GeoNames matching and hierarchical clustering to identify affected regions from textual sources. Their method achieves high alignment accuracy between extracted clusters and article scope. SeqIA~\cite{lopez2025} adopts a transformer-based approach using RoBERTa for location recognition together with dependency parsing to resolve ambiguous place names.

Although these approaches improve the extraction of geographic information, they remain constrained by document-level assumptions and limited spatial representation. Many geoparsing pipelines implicitly assume a single dominant impact location per document, reducing their ability to represent spatially distributed drought impacts across multiple regions.

\subsection{Drought Impact Extraction}

Alongside geographic information extraction, recent studies have increasingly focused on identifying and categorising drought impacts directly from textual data. This task generally involves classifying text fragments into predefined impact categories such as agriculture, forestry, water supply, or wildfires.

Earlier approaches relied primarily on conventional machine learning methods using manually engineered textual features. Sodoge et al.~\cite{sodoge2023}, for example, apply lasso-regularised logistic regression on tf-idf representations, reporting high classification performance across several drought impact categories.

More recent systems employ transformer-based language models for relevance detection and impact classification. SeqIA~\cite{lopez2025} combines Longformer architectures with RoBERTa classifiers to identify drought-related content and assign impact categories using supervised learning approaches.

% Compared to location extraction, impact-type classification generally exhibits stronger and more stable performance across studies. This is partly due to clearer semantic distinctions between impact categories and the availability of annotated training datasets.

\subsection{Large Language Models for Information Extraction in climate research}

The emergence of large language models (LLMs) has created new opportunities for information extraction in geospatial and environmental applications. Unlike traditional supervised NLP systems, LLMs can perform zero-shot (no train samples) or few-shot extraction without requiring large labelled datasets.

Recent work has applied prompt-based LLMs to drought impact extraction. For example, CienaLLM~\cite{ciena2025} extracts drought impacts and affected administrative regions from Spanish news articles, achieving competitive impact classification while highlighting remaining challenges in location extraction.

Compared to traditional pipeline-based approaches, LLMs can jointly extract locations, impact descriptions, severity indicators, and temporal references within a single inference process, enabling more flexible representations of drought impacts.



% \subsection{Current Limitations}

% Despite recent advances, the extraction of high-resolution drought impact information from textual sources remains limited by several structural challenges. Existing methods often struggle to represent multiple geographically distributed impacts within a single document. 

% The reliability of extracted datasets is also affected by inconsistencies in reporting practices, linguistic ambiguity, and uneven spatial coverage in news sources. These limitations constrain the construction of structured drought impact datasets suitable for predictive modelling and operational forecasting applications \cite{sodoge2023, lopez2025, ciena2025}.

